% One-page letter / statement template. Matches the CV letterhead.
\documentclass[11pt, letterpaper]{article}
\usepackage[T1]{fontenc}
\usepackage[utf8]{inputenc}
\usepackage{charter}
\usepackage[top=0.55in, bottom=0.45in, left=0.8in, right=0.8in]{geometry}
\usepackage[dvipsnames]{xcolor}
\definecolor{linkblue}{HTML}{1155CC}
\usepackage[colorlinks=true, urlcolor=linkblue, linkcolor=linkblue]{hyperref}
\usepackage{parskip}
\pagestyle{empty}
\setlength{\parindent}{0pt}
\linespread{1.0}

\newcommand{\letterhead}{%
  \begin{center}
    {\large\bfseries Rodrigo França Chaves}\\[2pt]
    {\small
      Chicago, IL \enspace$\cdot$\enspace
      \href{mailto:rchaves@uchicago.edu}{rchaves@uchicago.edu} \enspace$\cdot$\enspace
      +1 312 217 5526 \enspace$\cdot$\enspace
      \href{https://rodrigofch7.github.io}{rodrigofch7.github.io}
    }
  \end{center}
  \vspace{2pt}
  {\rule{\linewidth}{0.8pt}}
  \vspace{10pt}
}

\begin{document}
\letterhead

\textbf{Econic Partners} \\
Research Analyst, 2027 start \\
\textit{Based in Chicago; open to any US office.}

\vspace{8pt}

Dear Recruiting Team,

I am applying for the 2027 Research Analyst position. I wrote to you earlier this year after visiting Econic in March with the University of Chicago Harris School, and your team kindly replied that this application would open and encouraged me to apply then. I have been waiting for it since.

The visit is what moved this from interest to intent. What struck me was that the analytical questions Econic works on, pricing, diversion and switching, market definition, entry and merger effects, are the same questions I have been working on for the past three years, in a different setting.

The closest thing I have to your work is my first-authored article in \textit{Utilities Policy}, which asks whether transferring Brazil's water and wastewater systems to private operators changed health outcomes. It compares propensity-matched municipalities before and after transfer between 1998 and 2021 using staggered difference-in-differences, and finds significant reductions in diarrheal morbidity among children under five, no effect on diseases unrelated to sanitation, and results that vary considerably across municipalities. The underlying exercise, estimating what a change of ownership and control did to outcomes when treatment arrives at different times in different places, is one I would recognize in a merger retrospective.

Two other pieces of my work map onto Econic's more directly than their subject matter suggests. At the World Bank this summer I quantified how developers in Mumbai substitute between competing regulatory instruments in the market for additional floor area, which is a switching question with prices attached. Separately, I built a horizontal-equity audit of New York's property tax roll, in which the entire methodological problem was defining a defensible comparison set: a citywide benchmark is useless when a Manhattan cooperative and a Queens two-family fall into the same column, so I built peer groups on borough, building class, tax class, decade built, size, and value band, across 1.1 million properties. Choosing what counts as comparable, and being able to defend that choice under attack, is market definition in another setting.

I am comfortable at the scale this work involves, in Stata, Python, and R, and I take the communication side seriously: I spent two years teaching econometrics and GIS while working as a research assistant, and I build my analyses to be interrogated by people who cannot read the code.

I have not worked in economic consulting, and I would be starting rather than arriving. But I am completing my Master's in Computational Analysis and Public Policy in June 2027, which fits this cycle, and I would be joining a firm I have already visited and already want to work for.

Thank you for your consideration. I would welcome the chance to talk.

\vspace{6pt}

Sincerely, \\
Rodrigo França Chaves

\end{document}
